% !TEX root = ../../../main.tex
% Block card: catalogue.gramian_fusion -- "Gramian fusion" (Part I, the blocks chapter).
\cardfield{What it does}
Makes one colour picture from three: the summation field as red, the difference field as
green and the recurrence plot as blue (see the card of \blockref{catalogue.gramian_gasf}
for what each one shows).

\cardfield{When to use it}
To see on one span the picture the fusion networks were trained on. To build such pictures
for many windows, ready for \blockref{catalogue.cnn_score}, use
\blockref{catalogue.window_images}, which makes the same picture window by window.

\cardfield{What goes in, what comes out}
In goes a span of one channel, at most 5{,}000 samples. Out comes a square colour image with
one row and one column per sample.

\cardfield{Settings that matter}
\emph{m} and \emph{tau} shape the blue layer only, as in
\blockref{catalogue.gramian_recurrence}. Leave them at 3 and 4 if the picture is going to a
trained network: that is what the networks saw.

\cardfield{Watch out for}
Each colour is stretched to its own full range, so the colours say where a pixel sits within
its own layer, not how strong anything is. A single missing sample, or a perfectly flat
span, turns the whole picture black.
